% ============================================================
%  Наративний огляд літератури (вдосконалена версія)
%  Тема: Методика організації групової взаємодії учнів
%        у процесі пізнавальної діяльності
%        на заняттях «Я досліджую світ»
% ============================================================
%  Компілятор: XeLaTeX  |  Бібліографія: biber
%  Збірка:
%    xelatex literature_review_improved
%    biber   literature_review_improved
%    xelatex literature_review_improved
%    xelatex literature_review_improved
% ============================================================

\documentclass[a4paper,14pt]{extarticle}

% ── Кодування та мова ─────────────────────────────────────
\usepackage{fontspec}
\usepackage{polyglossia}
\setdefaultlanguage{ukrainian}
\setotherlanguages{english,russian}
% Times New Roman (доступний на Overleaf/Windows/macOS)
% Якщо Times New Roman відсутній на Linux-сервері — замініть на:
% \setmainfont{FreeSerif}
\setmainfont{Times New Roman}

\newfontfamily\cyrillicfont{Times New Roman}
\newfontfamily\cyrillicfontsf{Arial}
\newfontfamily\cyrillicfonttt{Courier New}

\setmonofont{Courier New}

% ── Поля сторінки (КДПУ: л=30, п=10, в=20, н=20 мм) ─────
\usepackage[
  top    = 20mm,
  bottom = 20mm,
  left   = 30mm,
  right  = 10mm
]{geometry}

% ── Міжрядковий інтервал 1,5 ─────────────────────────────
\usepackage{setspace}
\doublespacing

% ── Відступ першого рядка ─────────────────────────────────
\usepackage{indentfirst}
\setlength{\parindent}{1.25cm}

% ── Нумерація сторінок — верхній правий кут ───────────────
\usepackage{fancyhdr}
\pagestyle{fancy}
\fancyhf{}
\fancyhead[R]{\thepage}
\renewcommand{\headrulewidth}{0pt}
\fancypagestyle{plain}{%
  \fancyhf{}
  \fancyhead[R]{\thepage}
  \renewcommand{\headrulewidth}{0pt}
}

% ── Заголовки (ДСТУ-стиль) ────────────────────────────────
% \section → ВЕЛИКИМИ ЛІТЕРАМИ, по центру, жирний, без крапки
% \subsection → з абзацного відступу, перша літера велика, жирний
\usepackage{titlesec}
\titleformat{\section}[block]
  {\normalfont\normalsize\bfseries\centering\MakeUppercase}
  {\thesection}{1em}{}
\titleformat{\subsection}[block]
  {\normalfont\normalsize\bfseries}
  {\thesubsection}{1em}{}
\titlespacing*{\section}{0pt}{12pt}{6pt}
\titlespacing*{\subsection}{1.25cm}{10pt}{4pt}

% ── Зміст ─────────────────────────────────────────────────
\usepackage{tocloft}
\renewcommand{\cftsecfont}{\normalfont}
\renewcommand{\cftsecpagefont}{\normalfont}
\renewcommand{\cftsubsecfont}{\normalfont}
\setlength{\cftbeforesecskip}{4pt}
\setlength{\cftsecindent}{0pt}
\setlength{\cftsubsecindent}{1.25cm}

% ── Цитування (APA-подібний стиль) ───────────────────────
\usepackage[
  backend  = biber,
  style    = authoryear,
  sorting  = nyt,
  maxnames = 3,
  minnames = 1,
  uniquelist = false,
  dashed   = false
]{biblatex}
\usepackage{csquotes}
\addbibresource{references_improved.bib}

% ── Таблиці та малюнки ────────────────────────────────────
\usepackage{booktabs}
\usepackage{longtable}
\usepackage{array}
\usepackage{caption}
\usepackage{graphicx}
\captionsetup{font=normalsize, labelfont=bf, justification=centering,
              labelsep=period, skip=4pt}
% Нумерація таблиць/рисунків у межах розділу: Таблиця 2.1, Рис. 3.1
\usepackage{chngcntr}
\counterwithin{table}{section}
\counterwithin{figure}{section}
\renewcommand{\thetable}{\arabic{section}.\arabic{table}}
\renewcommand{\thefigure}{\arabic{section}.\arabic{figure}}

% ── Гіперпосилання ────────────────────────────────────────
\usepackage[
  colorlinks = true,
  linkcolor  = black,
  citecolor  = black,
  urlcolor   = blue,
  unicode    = true,
  pdfencoding=auto
]{hyperref}

% ── Додаткові пакети ──────────────────────────────────────
\usepackage{microtype}
\usepackage{enumitem}
\usepackage{xcolor}
\usepackage{multicol}

% ── Лапки ─────────────────────────────────────────────────
% «Ялинки» для українського тексту
\usepackage{upquote}

% ── Метадані PDF ──────────────────────────────────────────
\hypersetup{
  pdftitle  = {Методика організації групової взаємодії учнів на заняттях «Я досліджую світ»},
  pdfauthor = {Студент спеціальності 013 Початкова освіта},
  pdflang   = {uk}
}

% ============================================================
\begin{document}

% ============================================================
% ТИТУЛЬНИЙ АРКУШ
% ============================================================
\begin{titlepage}
\begin{center}
\MakeUppercase{Міністерство освіти і науки України}\\[0.3em]
\MakeUppercase{Криворізький державний педагогічний університет}\\[0.3em]
\MakeUppercase{Кафедра початкової освіти}

\vspace{1.5cm}
{\small Допущено до захисту}\\
{\small Завідувач кафедри \underline{\hspace{4cm}}}\\
{\small «\underline{\hspace{1cm}}» \underline{\hspace{3cm}} 20\underline{\hspace{0.7cm}} р.}

\vspace{2.5cm}
{\large\bfseries
Методика організації групової взаємодії учнів\\
у процесі пізнавальної діяльності\\
на заняттях «Я досліджую світ»}\\[0.5em]
{\normalsize Наративний огляд літератури}

\vspace{1.5cm}
{\normalsize Курсова робота\\
зі спеціальності 013 Початкова освіта}

\vspace{2cm}
\begin{flushright}
{\normalsize
Виконав(ла): студент(ка) \underline{\hspace{0.5cm}} курсу,\\
групи \underline{\hspace{2cm}}\\
\underline{\hspace{5cm}}\\
(Прізвище Ім'я По-батькові)\\[1em]
Науковий керівник:\\
\underline{\hspace{5cm}},\\
к.\,пед.\,н., доцент
}
\end{flushright}

\vfill
{\normalsize Кривий Ріг~--- 2025}
\end{center}
\end{titlepage}

\setcounter{page}{2}

% ============================================================
% ЗМІСТ
% ============================================================
\tableofcontents
\newpage

% ============================================================
% МЕТОДОЛОГІЧНА ПРИМІТКА
% ============================================================
\section*{Методологічна примітка}
\addcontentsline{toc}{section}{Методологічна примітка}

Відповідно до статті~8, частини~4 Закону України «Про освіту» (2017) та Кодексу академічної доброчесності Криворізького державного педагогічного університету, автор зазначає таке. У процесі підготовки цього огляду систему штучного інтелекту Claude (Anthropic) було залучено виключно для технічних завдань: (а)~первинної фільтрації та сортування бібліографічних записів із бази даних Scopus за заздалегідь визначеними критеріями релевантності; (б)~тематичного групування відібраних джерел відповідно до структури огляду; (в)~форматування BibTeX-файлу і верстки LaTeX-документа згідно з вимогами методичних рекомендацій. Усі аналітичні судження, критичне оцінювання змісту джерел, перефразування, синтез ідей, формулювання висновків та виявлення дослідницьких прогалин є результатом самостійної інтелектуальної роботи автора. Використання інструментів штучного інтелекту для зазначеної технічної підтримки не є порушенням академічної доброчесності за умови його прозорого розкриття, що цим і здійснюється.

\newpage

% ============================================================
% РОЗДІЛ 1. ВСТУП
% ============================================================
\section{Вступ}

Групова взаємодія учнів у пізнавальній діяльності є одним із ключових механізмів реалізації компетентнісного підходу в сучасній початковій школі. В умовах Нової української школи (НУШ) інтегрований курс «Я досліджую світ» (ЯДС) створює унікальний навчальний контекст, що органічно поєднує природничу, соціальну, громадянознавчу та здоров'язбережувальну освітні галузі в єдиному дослідницькому просторі. За типовою освітньою програмою для 3--4 класів (авт. О.\,Я.~Савченко), заняття ЯДС передбачають три форми учнівської дослідницької активності~--- дослідження-розпізнавання, дослідження-спостереження та дослідження-пошуку~--- реалізація яких у груповому форматі відповідає природній соціальній потребі молодших школярів у спільній діяльності \parencite{Savchenko2019}.

Незважаючи на значний масив міжнародних досліджень з кооперативного (cooperative learning) та дослідницького (inquiry-based learning, IBL) навчання, питання методики організації саме групової взаємодії в контексті інтегрованих курсів початкової школи залишається недостатньо систематизованим у вітчизняній педагогічній науці. Наявна суперечність між добре опрацьованою в зарубіжній педагогіці теорією групового навчання і недостатньою її адаптацією до специфіки НУШ зумовлює актуальність цього огляду.

\textbf{Мета} огляду~--- систематизувати міжнародні та вітчизняні наукові дані щодо психолого-педагогічних засад, типів, вікових особливостей та методичних підходів до організації групової взаємодії учнів початкової школи у процесі пізнавальної діяльності дослідницького характеру в контексті курсу ЯДС.

\textbf{Завдання огляду}: (1)~з'ясувати теоретичні засади кооперативного та дослідницького навчання; (2)~охарактеризувати типи групової роботи та вікові особливості їх організації; (3)~розкрити роль учителя-фасилітатора; (4)~систематизувати методичні підходи до структурування завдань та оцінювання; (5)~визначити специфіку інтеграції змістів у рамках курсу ЯДС; (6)~виявити прогалини в дослідженнях і накреслити напрями подальших студій.

\textbf{Методологія пошуку і відбору джерел.} Пошук здійснювався у базі даних Scopus за тринадцятьма тематичними запитами, об'єднаними у п'ять блоків (детально описано в Додатку~А). Хронологічний діапазон: 2015--2026 роки; тип документу: статті та огляди, мова~--- англійська і українська. Критерії включення: емпіричні або оглядові публікації, присвячені груповій/кооперативній/колаборативній взаємодії учнів початкової школи (вік 6--11 років) у пізнавальній або дослідницькій діяльності. Критерії виключення: дошкільна і старша школа, вища освіта, роботи без зв'язку з класно-урочним контекстом. Крім того, до огляду залучено ключові вітчизняні наукові праці та нормативні документи НУШ, відібрані за принципом авторитетності джерела. За результатами первинної фільтрації (280 записів) і відбору за критеріями релевантності до огляду включено 63~джерела, з яких безпосередньо процитовано 46.

\newpage

% ============================================================
% РОЗДІЛ 2. ТЕОРЕТИЧНІ ЗАСАДИ ГРУПОВОЇ ВЗАЄМОДІЇ
% ============================================================
\section{Теоретичні засади групової взаємодії учнів}

\subsection{Психолого-педагогічні основи кооперативного навчання}

Кооперативне навчання є однією з найбільш доказово обґрунтованих стратегій організації навчального процесу в початковій школі. Теоретичним підґрунтям цього підходу слугує розроблена Дж.~та Р.~Джонсонами модель позитивної взаємозалежності, де спільне розв'язання завдань активізує когнітивні й метакогнітивні процеси, недосяжні в умовах індивідуальної роботи: учні вербалізують міркування, отримують зворотний зв'язок від однолітків і коригують власні способи дій \parencite{Gu2015143}. Вітчизняна педагогічна наука розвиває цей підхід у концепції інтерактивного навчання О.\,І.~Пометун і Л.\,В.~Пироженко, де групова взаємодія розглядається як обов'язкова умова формування ключових компетентностей молодших школярів, а не лише засіб підвищення успішності \parencite{Pometun2004}.

Систематичний мета-аналіз, що охоплює тридцять емпіричних досліджень, підтвердив статистично значущий позитивний вплив кооперативного навчання на математичну успішність учнів початкової школи, наголошуючи водночас на тому, що результативність стратегії залежить від якості її методичного впровадження~--- зокрема від ступеня структурування взаємозалежності й чіткості розподілу ролей \parencite{Talkhan2025}. Квазіекспериментальне дослідження, проведене в початкових школах Італії за участі тридцяти шести класів, показало, що структуроване кооперативне навчання із технологічною підтримкою суттєво розвиває соціальні й емоційні навички молодших школярів порівняно з контрольною групою \parencite{Zagni2025}. Довгострокове лонгітюдне спостереження учнів п'ятого і шостого класів підтвердило, що кооперативна методологія асоціюється зі стабільнішим розвитком творчого мислення та лексичного збагачення протягом усього пізнього дитинства \parencite{Segundo-Marcos2025}.

Ефективне впровадження кооперативного навчання потребує цілеспрямованої підготовки вчителя і врахування специфіки класного середовища. Якісне дослідження на основі глибинних інтерв'ю з тридцятьма двома вчителями початкової школи в Ірані виявило три ключові групи перешкод: організаційні (велика наповнюваність класів, брак часу), педагогічні (нестача методичної компетентності) та пов'язані з учнівською поведінкою (різний рівень підготовки, конфлікти в групі) \parencite{Keramati2025120}. Ці перешкоди є характерними і для українського шкільного контексту~--- особливо в умовах змішаних класів і значної диференціації рівнів навчальних досягнень, типових для НУШ після переходу з дистанційного формату. Пілотна програма у гетерогенних класах Женеви продемонструвала, що правильно організована кооперативна діяльність захищає соціальний клімат класу і знижує рівень ізоляції окремих учнів \parencite{Buchs202568}. Водночас аналіз динаміки групової колаборації у четвертому класі виявив, що позаурочні взаємодії між учнями виконують регулятивну функцію в структурі групової роботи, перерозподіляючи владні відносини всередині команди \parencite{Langer-Osuna2020514}. Ці дані вказують на те, що ефективність кооперативного навчання не є автоматичним наслідком групування учнів, а залежить від усвідомленого педагогічного управління груповою взаємодією~--- висновок, що прямо кореспондує з вимогами Державного стандарту НУШ щодо компетентнісно-орієнтованого навчання \parencite{Bibik2017}.

\subsection{Типи групової роботи та їх порівняльна характеристика}

У сучасній педагогіці розрізняють кілька форм організації спільної навчальної діяльності, що відрізняються за ступенем взаємозалежності учасників, характером завдання і рівнем автономії групи. Систематизація цих форм є необхідною умовою усвідомленого педагогічного проєктування, оскільки їх підміна або змішування знижує ефективність навчання.

\textit{Кооперативне навчання} (cooperative learning) передбачає чіткий структурований розподіл ролей, позитивну взаємозалежність результатів та індивідуальну відповідальність кожного члена групи за кінцевий спільний продукт. \textit{Колаборативне навчання} (collaborative learning) є більш гнучкою, рівноправною формою взаємодії, де акцент робиться на спільному конструюванні смислів та відкритому обміні ідеями без жорсткого розподілу функцій. \textit{Проєктна робота} (project-based learning, PjBL) поєднує ознаки обох форм, але орієнтується на виготовлення конкретного продукту і вирішення автентичної проблеми. \textit{Дослідницька групова діяльність} у форматі IBL спрямована передусім на колективне породження та перевірку гіпотез, тобто акцентує процес пізнання, а не продукт.

\begin{center}
\captionof{table}{Порівняльна характеристика форм групової навчальної діяльності}
\label{tab:2.1}
\begin{longtable}{|>{\raggedright\arraybackslash}p{2.8cm}|
                    >{\raggedright\arraybackslash}p{3cm}|
                    >{\raggedright\arraybackslash}p{3.2cm}|
                    >{\raggedright\arraybackslash}p{3.2cm}|}
\hline
\textbf{Характеристика} &
\textbf{Кооперативне навчання} &
\textbf{Колаборативне навчання} &
\textbf{Дослідницька (IBL) групова діяльність} \\
\hline
Розподіл ролей & Чіткий, структурований & Гнучкий, спільно визначений & Ситуативний, за потребою \\
\hline
Взаємозалежність & Позитивна, запрограмована & Природна, виникає органічно & Спільна навколо наукової проблеми \\
\hline
Відповідальність & Індивідуальна + групова & Переважно групова & Групова, процес-орієнтована \\
\hline
Кінцевий результат & Спільний продукт & Спільне розуміння & Відповідь/пояснення явища \\
\hline
Роль учителя & Організатор, контролер & Фасилітатор & Тьютор, «провокатор» \\
\hline
Ступінь відкри\-тості завдання & Низький~--- середній & Середній & Середній~--- високий \\
\hline
Вік (оптимально) & 1–4 класи & 3–4 класи & 2–4 класи \\
\hline
Зв'язок з ЯДС & Структуровані дослідження-розпізнавання & Обговорення та побудова концептів & Дослідження-спостереження, дослідження-пошуку \\
\hline
\end{longtable}
\end{center}

Дослідження поведінки першокласників (7 років) у малих групах за методологією Lesson Study засвідчило, що навіть наймолодші учні здатні до продуктивної кооперації за умови відповідного педагогічного супроводу, спростовуючи поширене переконання про «незрілість» дітей цього віку для групової роботи \parencite{Jakavonytė-Staškuvienė2025}. Водночас дослідження спільного сторітелінгу в учнів 1--3 класів (97 осіб) показало, що тангібельні цифрові форми колаборації стимулюють більш інклюзивні та просоціальні взаємодії порівняно з нецифровими форматами \parencite{Filosofi2025801}. Аналіз педагогіко-технологічних характеристик платформ електронного навчання виявив, що найефективніші підтримують усі фази проєктної роботи \parencite{Baziukė2025}. Дослідницька групова діяльність у природничій освіті початкової школи набуває форми запитання-орієнтованого навчання \parencite{Sakamoto2016105}. Вибір конкретного типу групової роботи для занять ЯДС слід здійснювати з урахуванням тематичного блоку програми, вікових особливостей учнів і наявних ресурсів: наприклад, кооперативні структури більш придатні для вивчення нового матеріалу, тоді як дослідницькі~--- для закріплення і поглиблення знань через самостійний пошук.

\subsection{Вікові особливості групової взаємодії молодших школярів}

Молодший шкільний вік охоплює діапазон від 6--7 до 10--11 років і характеризується інтенсивним розвитком довільності, рефлексії та здатності координувати власні дії з діями партнерів. Проте цей вік не є однорідним: між першокласниками (6--7 років) і четвертокласниками (9--10 років) існують суттєві відмінності в соціально-когнітивному розвитку, що безпосередньо впливає на характер групової взаємодії.

Учні 1--2 класів (6--8 років) потребують максимального структурування групових завдань, чіткого розподілу ролей і постійної педагогічної підтримки. Поширене переконання про те, що першокласники ще не здатні до продуктивної групової роботи, спростовується емпіричними даними: цілеспрямоване дослідження поведінки семирічних учнів у малих групах із трьох осіб за методологією Lesson Study показало, що за відповідного педагогічного супроводу учні першого класу виявляють стійку здатність до кооперативної взаємодії \parencite{Jakavonytė-Staškuvienė2025}. Учні другого класу здатні виконувати inquiry-завдання з природознавства і технологій, розвиваючи при цьому наукову грамотність у суміщеному дослідницько-груповому форматі \parencite{Zuñiga-Quispe2025365}.

Учні 3--4 класів (8--11 років), на яких зосереджений курс ЯДС, демонструють вищу когнітивну гнучкість, здатність до аргументації та рефлексії власних дій, що відкриває можливості для більш відкритих і проблемних форм групової діяльності. На матеріалі четвертокласників встановлено нелінійний характер зв'язку між ступенем відкритості дослідницького завдання і рівнем навчальної автономії: широкий простір для IBL одночасно провокує як найвищі, так і найнижчі рівні самостійності, що вимагає диференційованого підходу з боку вчителя \parencite{Aminga2026}. Типова освітня програма ЯДС для 3--4 класів враховує цю специфіку, орієнтуючи педагогів на поступове розширення ступеня самостійності учнів у дослідницьких активностях~--- від дослідження-розпізнавання до дослідження-пошуку \parencite{Savchenko2019}.

Соціально-емоційні умови групової взаємодії набувають особливого значення у молодшому шкільному віці. Дослідження управління «фейсворком» (facework)~--- підтримання соціального іміджу і гідності учасників взаємодії (концепція Е.~Гоффмана)~--- виявило, що позитивна соціальна атмосфера у групі є необхідною умовою продуктивного пізнавального обміну між однолітками, тоді як загроза публічної невдачі суттєво пригнічує навчальний діалог \parencite{Evnitskaya2021169}. Аналіз групової динаміки четвертокласників під час колаборативного розв'язання математичних задач підтверджує, що відносини влади й статусу всередині групи є рухомими і потребують усвідомленого педагогічного управління \parencite{Langer-Osuna2020514}. Отже, ефективна організація групової взаємодії молодших школярів вимагає одночасного врахування як вікових когнітивних можливостей, так і соціально-емоційного контексту спільної діяльності.

\subsection{Концептуальна модель організації групової дослідницької взаємодії}

На основі синтезу розглянутих теоретичних підходів можна запропонувати концептуальну модель організації групової дослідницької взаємодії учнів 3--4 класів у рамках курсу ЯДС (рис.~\ref{fig:2.1}). Модель охоплює чотири взаємопов'язані блоки: \textit{умови} (інтегрований зміст ЯДС, вікові характеристики, матеріальне середовище), \textit{процеси взаємодії} (кооперативні, колаборативні, дослідницькі), \textit{педагогічний супровід} (рівень фасилітації, тип завдання, оцінювання) та \textit{результати} (предметні, метапредметні, соціально-емоційні). Стрілки двосторонні~--- модель відображає циклічний, а не лінійний характер пізнавальної діяльності. Кожен з цих блоків розгорнуто у наступних розділах огляду.

\begin{figure}[h!]
\centering
\setlength{\unitlength}{1cm}
\begin{picture}(14,8)
% Блок 1: Умови
\put(0,5.5){\framebox(3.8,2){\shortstack{\textbf{УМОВИ}\\ \footnotesize Зміст ЯДС\\ \footnotesize Вікові особливості\\ \footnotesize Середовище}}}
% Блок 2: Процеси
\put(5,5.5){\framebox(4,2){\shortstack{\textbf{ПРОЦЕСИ}\\ \footnotesize Кооперація\\ \footnotesize Колаборація\\ \footnotesize IBL-діяльність}}}
% Блок 3: Педагогічний супровід
\put(0,0.5){\framebox(3.8,2){\shortstack{\textbf{ПЕДАГОГІЧНИЙ}\\ \textbf{СУПРОВІД}\\ \footnotesize Тип завдання\\ \footnotesize Фасилітація\\ \footnotesize Оцінювання}}}
% Блок 4: Результати
\put(5,0.5){\framebox(4,2){\shortstack{\textbf{РЕЗУЛЬТАТИ}\\ \footnotesize Предметні\\ \footnotesize Метапредметні\\ \footnotesize Соціально-емоційні}}}
% Зв'язки
\put(3.8,6.5){\vector(1,0){1.2}}
\put(5,6.5){\vector(-1,0){1.2}}
\put(3.8,1.5){\vector(1,0){1.2}}
\put(5,1.5){\vector(-1,0){1.2}}
\put(2,5.5){\vector(0,-1){1.2}}
\put(2,3.7){\vector(0,1){1.2}}
\put(7,5.5){\vector(0,-1){1.2}}
\put(7,3.7){\vector(0,1){1.2}}
\end{picture}
\caption{Концептуальна модель організації групової дослідницької взаємодії учнів на заняттях ЯДС}
\label{fig:2.1}
\end{figure}

\newpage

% ============================================================
% РОЗДІЛ 3. ПІЗНАВАЛЬНА ДІЯЛЬНІСТЬ У ГРУПОВОМУ ФОРМАТІ
% ============================================================
\section{Пізнавальна діяльність у груповому форматі}

\subsection{Дослідницька поведінка учнів та inquiry-based learning}

Дослідницьке навчання (inquiry-based learning, IBL) як педагогічна стратегія ґрунтується на активному залученні учнів до процесів спостереження, формулювання запитань, висування гіпотез, збирання даних та осмислення результатів. Квазіекспериментальне дослідження ефективності IBL у змішаному синхронному навчальному середовищі підтвердило, що цей підхід суттєво підвищує залученість і навчальні результати учнів початкової школи порівняно з традиційним навчанням \parencite{Xu2026}. Якісний мета-аналіз міжнародних досліджень IBL засвідчив, що дослідницькі моделі навчання одночасно розвивають соціальні навички та покращують академічну успішність молодших школярів у різних культурних контекстах \parencite{Seprie2025}. Важливо, що ці ефекти не є рівномірними: відкриті форми IBL дають найкращі результати у класах з достатнім рівнем попередньої академічної підготовки, тоді як у класах із суттєвим рівневим розривом ефективнішими виявляються структуровані та спрямовані форми IBL \parencite{Aminga2026}.

Дослідницька поведінка учнів початкової школи набуває різних форм залежно від навчального предмета і характеру завдання. Дослідження формування творчого мислення в учнів початкової школи на матеріалі природознавства підтвердило переваги дослідницько-орієнтованого підходу перед традиційним викладом \parencite{Leasa2025318}. Систематизація досвіду викладання природознавства і технологій у другому класі демонструє, що IBL-підхід ефективно розвиває наукову грамотність вже у ранньому шкільному віці \parencite{Zuñiga-Quispe2025365}. Модель дослідницького навчання «Сім'я–Школа–Суспільство», розроблена в рамках китайської освітньої реформи, ілюструє потенціал розширення дослідницького середовища учнів за межі класної кімнати \parencite{Xia2025340}. Взаємодія учнів з фізичними симуляціями у форматі IBL сприяє осмисленню абстрактних понять через активне маніпулювання і перевірку гіпотез; порівняння вчительської та учнівської моделей взаємодії з симуляціями виявило різні профілі пізнавальної активності \parencite{Magkouta2024253}.

Окремий вектор досліджень пов'язаний із IBL у контексті екологічної освіти та громадянської відповідальності~--- аспекту, прямо передбаченого змістом курсу ЯДС. Якісне дослідження, присвячене формуванню «екологічного громадянства» у молодших школярів, показало, що inquiry-практики, спрямовані на вивчення місцевого біорізноманіття, розвивають як природничі компетентності, так і ціннісні орієнтації учнів щодо навколишнього середовища \parencite{Christodoulou2025969}. Тривале кейс-дослідження (13 місяців) засвідчило, що систематична дослідницька діяльність учнів 4--5 класів, опосередкована технологічним середовищем, формує стійкі пізнавальні стратегії та навички спільного конструювання знань \parencite{Viilo2018407}. Наведені дані дозволяють розглядати IBL як методологічну основу пізнавальної діяльності учнів на заняттях ЯДС, проте з урахуванням необхідності диференціації ступеня відкритості завдань залежно від рівня готовності класу.

\subsection{Роль учителя-фасилітатора}

Перехід від моделі вчителя-транслятора знань до моделі вчителя-фасилітатора є однією з ключових характеристик сучасної педагогіки початкової школи, зокрема в контексті IBL і кооперативного навчання. Вітчизняна наукова традиція пов'язує цей перехід із концепцією суб'єкт-суб'єктної взаємодії, розробленою у працях В.\,Ф.~Паламарчук, де вчитель виступає не джерелом готових знань, а організатором самостійного пошуку та «диригентом» думок \parencite{Palamarchuk2000}. Порівняльне дослідження практик двох учителів початкової школи Південної Африки з різним педагогічним стажем показало, що досвідчені педагоги більш гнучко чергують пряму інструкцію і відкриті дослідницькі ситуації, тоді як початківці схильні до надмірного структурування завдань~--- патерн, характерний і для молодих педагогів НУШ \parencite{Photo20251493}. Аналіз даних австрійського TIMSS-2019 дослідження виявив, що вчителі початкової школи, які частіше застосовують дослідницькі стратегії навчання, демонструють вищу чутливість до складу класу і гнибкіше адаптують інструкцію до потреб учнів \parencite{Neubauer-Hametner2025}.

Дослідження виокремлюють кілька рівнів фасилітативного втручання, що відповідають різним типам групової дослідницької активності. Кейс-дослідження у чотирьох норвезьких початкових класах засвідчило, що перехід від закритих питань і оцінювальних реакцій до відкритих запитань і підтримки учнівських ідей суттєво підвищує рівень спільного конструювання знань у класній дискусії \parencite{Hogsnes20241}. Цей висновок узгоджується з ідеєю «діалогічного навчання» (dialogic teaching), де педагог не оцінює відповіді учнів, а будує на них подальший дослідницький діалог. Однорічна програма професійного розвитку підтвердила, що вчителі початкової школи здатні опанувати фасилітацію академічної аргументації та діалогічного обговорення текстів за умови систематичної методичної підтримки \parencite{Reznitskaya2019254}. Цей висновок є принциповим для НУШ: готовність педагогів до фасилітаторської ролі формується не стихійно, а через цілеспрямовану систему неперервної освіти.

Технологічно опосередкована фасилітація відкриває нові можливості для організації IBL-навчання. Експериментальне дослідження показало, що мобільні підказки-орієнтири під час дослідницьких завдань поза класною кімнатою підвищують якість дослідницького процесу і знижують когнітивне навантаження на учнів \parencite{Martin20245}. Водночас дослідження досвіду студентів педагогічних спеціальностей із inquiry-послідовностями для початкової школи виявляє, що майбутні вчителі часто недооцінюють складність відкритих форм IBL~--- явище, документоване і серед практикуючих педагогів в Україні \parencite{Haslekås2022}. Виклики, пов'язані з роллю фасилітатора, включають: ризик «псевдо-IBL» (коли вчитель формально організовує дослідження, але фактично нав'язує відповідь), труднощі управління часом у великому класі та складність одночасного моніторингу кількох малих груп.

Для контексту ЯДС роль учителя-фасилітатора набуває додаткових вимірів. По-перше, інтегрований характер курсу передбачає, що педагог повинен рівною мірою орієнтуватися у природничому, соціальному та громадянознавчому змістах. По-друге, принцип «наскрізних умінь», закладений у державному стандарті, вимагає від учителя не просто допомагати учням знаходити відповіді, але й формувати культуру наукового запитування і спільного аналізу \parencite{Bibik2017}.

\newpage

% ============================================================
% РОЗДІЛ 4. МЕТОДИЧНІ ПІДХОДИ
% ============================================================
\section{Методичні підходи до організації групової роботи}

\subsection{Структурування завдань для групової дослідницької діяльності}

Якість групової дослідницької діяльності значною мірою визначається тим, як сформульоване і структуроване навчальне завдання. Концептуальна рамка проєктно-орієнтованого STEM-навчання (PjBL-STEM), розроблена для початкової школи, виокремлює кілька обов'язкових компонентів ефективного групового завдання: автентичність проблеми, міждисциплінарний характер, чіткий кінцевий продукт і вбудований механізм зворотного зв'язку \parencite{Retno2025579}. Реалізація цих принципів у практиці STEAM-проєктів початкової школи демонструє, що завдання, засновані на реальних суспільних проблемах (зокрема, кліматичних змінах і сталому розвитку), підвищують мотивацію учнів і забезпечують природну інтеграцію предметних змістів \parencite{Vicente20201}.

Типова освітня програма ЯДС для 3--4 класів окреслює тематичні модулі («Природа і я», «Людина серед людей», «Я~--- громадянин», «Я і здоров'я»), кожен із яких містить дослідницькі завдання, придатні для групового виконання \parencite{Savchenko2019}. Структурування цих завдань відповідно до рівня відкритості IBL-циклу є методичним завданням педагога: завдання першого рівня (структуроване IBL) задають мету, метод і очікуваний результат; другого рівня (спрямоване IBL) задають проблему і мету, але залишають метод відкритим; третього рівня (відкрите IBL) передбачають самостійний вибір дослідницького питання і методу \parencite{Aminga2026}.

Дослідження реалізації трьох дизайн-орієнтованих STEM-проєктів у початкових класах показало, що завдання з елементами інженерного проєктування суттєво розвивають наукову творчість і вимагають від групи узгодженого переходу між фазами ідеації, прототипування й оцінювання \parencite{Zhang2025568}. Розробка та апробація інтегрованих засобів навчання природознавства і суспільствознавства, орієнтованих на метакогнітивний розвиток, засвідчила, що тематичне структурування міждисциплінарних завдань покращує критичне мислення молодших школярів \parencite{Sururuddin2026}. Аналіз кооперативного викладання STEAM у початковій школі за design-based підходом підтверджує, що спільна розробка і виконання проєктних завдань є ефективним механізмом інтеграції предметних знань \parencite{Li2022}.

Цифрові платформи відкривають нові можливості для структурування і підтримки групової дослідницької діяльності. Оцінювання педагогіко-технологічних характеристик платформ виявило, що найефективніші інструменти підтримують кілька фаз спільної роботи~--- від постановки завдання і розподілу ролей до документування дослідницького процесу і презентації результатів \parencite{Baziukė2025}. Проте необхідно враховувати нерівномірність технічного оснащення шкіл в Україні та різний рівень цифрової компетентності вчителів: у 2023--2024 роках понад 40\% початкових класів здійснювали навчання в умовах обмеженого або нестабільного доступу до інтернету, що суттєво впливає на вибір цифрових форматів підтримки групової роботи.

\subsection{Оцінювання групової взаємодії}

Оцінювання в контексті групової дослідницької діяльності є методично складним завданням, що потребує поєднання індивідуальних та групових показників, врахування як процесу, так і продукту спільної роботи. Принципова відмінність між оцінюванням \textit{продукту} (що група зробила) і оцінюванням \textit{процесу} (як відбувалася взаємодія) є вихідною позицією сучасних підходів до оцінювання груповий роботи в початковій школі \parencite{Pometun2004}.

Дослідження впливу Програми оцінювання кооперативного навчання на соціальну взаємодію учнів 3--6 класів початкової школи Китаю (165 учасників) показало, що поєднання індивідуальної і групової складових оцінювання і фокус одночасно на академічних та соціальних результатах суттєво підвищує якість міжособистісної взаємодії в групах \parencite{Zeng2025317}. Ці дані узгоджуються з результатами дослідження, де залучення учнів до розробки критеріїв взаємооцінювання за допомогою мобільних пристроїв позитивно вплинуло на їхні досягнення та усвідомленість власних навчальних стратегій \parencite{Lai2015149}.

Взаємооцінювання є особливо ефективним у контексті дослідницького розв'язання задач. Якісний аналіз процесу розв'язання математичних задач шестикласниками засвідчив, що взаємооцінювання як складова нелінійного орієнтувального базису суттєво прискорює еволюцію дослідницьких стратегій групи і підвищує усвідомленість кожного учасника щодо якості власних міркувань \parencite{Torregrosa202189}. Рамка ECEE (дослідження~-- створення~-- вираження~-- оцінювання), розроблена для підготовки вчителів початкової школи, виокремлює оцінювання як окрему і рівнозначну фазу творчо-колаборативної діяльності, а не лише підсумковий захід \parencite{Mariati2026}.

Для молодших школярів особливу роль відіграють рубрики оцінювання~--- чіткі й доступні критерії, що одночасно орієнтують учнів у вимогах і стають інструментом самооцінювання та взаємооцінювання. Дослідження здатності майбутніх учителів проєктувати дослідницькі завдання виявило систематичні прогалини у формулюванні вимірюваних критеріїв оцінювання~--- насамперед у ситуаціях із несподіваними результатами \parencite{Lefkos2025349}. Стосовно курсу ЯДС рубрики оцінювання мають охоплювати щонайменше три виміри: (1)~якість дослідницького процесу (формулювання запитань, збір даних, перевірка гіпотез); (2)~рівень кооперативної взаємодії (розподіл ролей, комунікація, спільне прийняття рішень); (3)~презентацію та обговорення результатів. Оцінювання групових навчальних продуктів на цифрових навчальних платформах потребує, крім того, врахування колаборативних показників~--- ступеня залучення кожного члена групи до спільного продукту \parencite{Baziukė2025}.

\subsection{Групова взаємодія в умовах гібридного та дистанційного навчання}

Специфічний виклик для організації групової взаємодії в українських школах пов'язаний із реаліями воєнного часу: починаючи з 2022 року значна частина шкіл здійснює навчання у змішаному (гібридному) або повністю дистанційному форматі. Порівняльне дослідження ефектів кооперативного навчання в синхронному онлайн-форматі та в класі підтвердило, що обидва формати є результативними, однак відрізняються за характером міжособистісної взаємодії та груповою динамікою: онлайн-формат ускладнює невербальну комунікацію та спонтанні «перемовини» між учасниками групи \parencite{Ralević2025}. Тривале кейс-дослідження засвідчило, що систематична дослідницька діяльність, опосередкована технологічним середовищем, здатна формувати стійкі пізнавальні стратегії навіть за умов дистанційної роботи \parencite{Viilo2018407}.

Для занять ЯДС у гібридному форматі ключовими є такі методичні рішення: (а)~забезпечення синхронного спільного доступу до дослідницьких матеріалів (цифрові науково-дослідницькі щоденники, спільні презентації); (б)~адаптація фізичних дослідів до формату домашнього дослідження; (в)~використання асинхронних каналів зворотного зв'язку між учасниками групи. Дослідження показало, що мобільні підказки-орієнтири під час позакласних IBL-активностей здатні ефективно підтримувати дослідницький процес і поза школою \parencite{Martin20245}. Аналіз сприйняття кооперативного навчання учнями сільських початкових шкіл Китаю зафіксував, що більшість учнів позитивно оцінює групову роботу навіть в умовах обмеженого доступу до ресурсів \parencite{Xiao2025117}~--- позитивний сигнал для реалізації ЯДС в українських школах різного типу.

\newpage

% ============================================================
% РОЗДІЛ 5. ІНТЕГРАЦІЯ ЗМІСТІВ У КУРСІ «Я ДОСЛІДЖУЮ СВІТ»
% ============================================================
\section{Інтеграція природничих, соціальних та громадянознавчих змістів
у курсі «Я досліджую світ»}

\subsection{Теоретичні засади міждисциплінарної інтеграції}

Інтегрований курс ЯДС у рамках НУШ поєднує природничу, соціальну, громадянознавчу та здоров'язбережувальну освітні галузі в єдиний навчальний простір, орієнтований на дослідницьку діяльність учнів. Типова освітня програма 3--4 класів (авт. О.\,Я.~Савченко) передбачає використання інтегрованих завдань дослідницького характеру, що охоплюють довкілля, суспільство та людину як взаємопов'язані об'єкти пізнання \parencite{Savchenko2019}. Теоретичним обґрунтуванням такого підходу є ідея про те, що розподіл знань на ізольовані дисципліни суперечить природі дитячого мислення~--- думка, що знаходить підтвердження у сучасних нейропедагогічних дослідженнях~--- тоді як інтеграція змістів дозволяє учням встановлювати міжпредметні зв'язки і бачити цілісність навколишнього світу \parencite{Lehane20261}.

Досвід реалізації інтегрованих підходів у початковій освіті різних країн підтверджує педагогічну обґрунтованість такого проєктування навчального змісту. Розробка і апробація тематичних засобів навчання, що інтегрують природничий і суспільствознавчий зміст у єдині дослідницькі завдання, суттєво покращує критичне мислення молодших школярів \parencite{Sururuddin2026}. Дослідження сприйняття вчителями початкової школи педагогіки глибинного навчання в умовах культурно-інтегрованої STEM-освіти засвідчило готовність педагогів до міждисциплінарного проєктування за наявності відповідної методичної підтримки \parencite{Putri20251331}. Важливо, що ефективна інтеграція не передбачає механічного поєднання змісту різних предметів: вона вимагає виявлення суттєвих зв'язків між поняттями та конструювання завдань, які ці зв'язки актуалізують.

STEM/STEAM-освіта є найбільш операціоналізованою формою інтеграції природничого та технологічного змістів на рівні початкової школи. Реалізація STEAM-проєктів із використанням освітньої роботики у п'ятому класі демонструє, як кооперативне навчання і проєктна діяльність органічно поєднуються в єдиному міждисциплінарному форматі \parencite{Vicente20201}. Design-based дослідження впровадження STEAM-навчання через кооперативне викладання підтвердило, що міждисциплінарне конструювання знань формує в учнів початкової школи інформаційну грамотність і дослідницькі компетентності вищого порядку \parencite{Li2022}. Практичні дослідницькі активності у природознавстві суттєво підвищують інтерес молодших школярів до STEM-дисциплін, формуючи стійку мотивацію до наукового пізнання вже у початкових класах \parencite{Ribarić2025}.

\subsection{Специфіка впровадження групових дослідницьких практик у курсі ЯДС}

Перенесення міжнародного досвіду групової IBL-діяльності на заняття ЯДС вимагає урахування кількох специфічних чинників. По-перше, тематична структура програми ЯДС охоплює широкий спектр об'єктів дослідження~--- від природних явищ («Зміни в природі», «Живі організми нашого краю») до соціальних реалій («Моя громада», «Я~--- громадянин України»). Ця тематична різноманітність вимагає від учителя гнучкого поєднання різних типів групових завдань: дослідження природних об'єктів органічно реалізується через структуровані та спрямовані форми IBL, тоді як вивчення суспільних тем~--- через колаборативне обговорення та проєктну діяльність. По-друге, інтегративний характер курсу відкриває можливість для «наскрізних» дослідницьких проєктів, що пов'язують природничий і суспільствознавчий зміст: наприклад, дослідження місцевого водойму (природничий аспект) може органічно включати вивчення громадянської відповідальності за довкілля (громадянознавчий аспект) \parencite{Christodoulou2025969}.

По-третє, важливою особливістю є облік регіональної та культурної специфіки. Типова освітня програма ЯДС передбачає широке залучення місцевого природного і соціального середовища як навчального ресурсу. Це означає, що групові дослідницькі завдання для учнів різних регіонів України будуть суттєво відрізнятися за своїм конкретним наповненням при збереженні спільної методологічної основи. Застосування ICT-середовищ із високим рівнем візуалізації для вивчення STEM-понять у початковій школі знижує рівень абстракції навчального матеріалу і стимулює дослідницьку допитливість учнів \parencite{Fanchamps2023102}~--- інструмент, що є особливо цінним в умовах дистанційного або гібридного навчання ЯДС.

По-четверте, необхідно враховувати суттєві виклики, задокументовані в дослідженнях: ризик «групових ефектів» (коли активні учні домінують, а пасивні «їдуть»); складність підтримання дослідницького темпу при великій наповнюваності класів; брак методичної підтримки для вчителів в умовах самостійного впровадження нових форм роботи. Вирішення цих проблем вимагає системних заходів~--- від розробки готових методичних сценаріїв групових занять ЯДС до вбудованих механізмів рефлексії (наприклад, щоденників дослідника або карток самооцінювання).

\newpage

% ============================================================
% ВИСНОВКИ
% ============================================================
\section*{Висновки огляду}
\addcontentsline{toc}{section}{Висновки огляду}

Проведений наративний огляд дозволяє сформулювати такі основні висновки.

\textit{По-перше}, кооперативне навчання в початковій школі є доказово ефективною стратегією, що одночасно розвиває академічні результати, соціально-емоційні навички та пізнавальну самостійність учнів, починаючи вже з першого класу. Різні форми групової роботи~--- кооперативна, колаборативна, проєктна та дослідницька~--- мають відмінні педагогічні профілі і мають добиратися відповідно до дидактичної мети та тематичного блоку програми ЯДС.

\textit{По-друге}, inquiry-based learning забезпечує методологічну основу для організації дослідницької пізнавальної діяльності: відкритість завдань, автентичність проблем і залучення дітей до наукових практик є ключовими чинниками ефективності цього підходу в початкових класах. Проте ступінь відкритості IBL-завдань має бути диференційованим залежно від рівня готовності класу та тематики: учні 3--4 класів здатні опрацьовувати завдання другого і третього рівня відкритості, якщо їм забезпечено відповідний педагогічний супровід.

\textit{По-третє}, роль учителя-фасилітатора є визначальною для якості групової взаємодії. Комунікативні стратегії педагога, рівень методичної готовності та здатність адаптувати ступінь підтримки до потреб учнів суттєво впливають на результативність спільної роботи. Перехід до фасилітаторської ролі є складним і тривалим процесом, що потребує цілеспрямованої підготовки у системі педагогічної освіти і неперервного розвитку в процесі практики.

\textit{По-четверте}, інтеграція природничих, соціальних та громадянознавчих змістів у дусі програми ЯДС є педагогічно обґрунтованою і потребує методично свідомого проєктування завдань, що актуалізують міжпредметні зв'язки у єдиній дослідницькій ситуації.

\textit{По-п'яте}, організація групової взаємодії в умовах гібридного та дистанційного навчання є самостійним методичним викликом, що потребує адаптації фізичних дослідів, цифрового підтримання взаємозалежності учасників групи та гнучкого підходу до оцінювання.

\textbf{Виявлені прогалини в дослідженнях}: (1)~майже відсутні дослідження, що безпосередньо вивчають групову взаємодію в умовах інтегрованих курсів типу ЯДС; (2)~недостатнє охоплення вікової підгрупи 3--4 класів як самостійного об'єкта вивчення~--- більшість досліджень оперують широким поняттям «початкова школа»; (3)~брак досліджень у контексті НУШ та пострадянських освітніх систем з урахуванням специфіки воєнного й поствоєнного часу; (4)~нестача емпіричних даних про ефективність різних форм оцінювання групової взаємодії саме у 3--4 класах з точки зору формувального підходу.

\textbf{Практичні рекомендації для вчителів}: розпочинати з кооперативних структур з чітким розподілом ролей, поступово переходячи до більш відкритих форм IBL у міру зростання рівня довіри і навичок групової роботи учнів; використовувати рубрики оцінювання, що враховують і процес взаємодії, і якість дослідницького продукту; регулярно рефлексувати власну комунікативну поведінку в ролі фасилітатора; при плануванні занять ЯДС визначати конкретний тип групової роботи (кооперативна, колаборативна, IBL) і добирати відповідний тип завдань.

Ці прогалини визначають актуальність і новизну подальшого дослідження в основній частині курсової роботи.

\newpage

% ============================================================
% СПИСОК ВИКОРИСТАНИХ ДЖЕРЕЛ
% ============================================================
\defbibnote{sourcesnote}{Джерела подані мовою оригіналу. Кириличні записи~--- спочатку,
           латинські~--- після.}


\printbibliography[
  title={Список використаних джерел},
  prenote=sourcesnote
]

\newpage

% ============================================================
% ДОДАТОК А — ПОШУКОВІ ЗАПИТИ
% ============================================================
\appendix
\section*{Додаток А}
\addcontentsline{toc}{section}{Додаток А. Пошукові запити в базі Scopus}

\begin{center}
\textbf{Пошукові запити, використані при підготовці огляду (база Scopus, 2015--2026)}
\end{center}

\begin{center}
\captionof{table}{Пошукові запити та результати первинної фільтрації}
\label{tab:A.1}
\begin{longtable}{|c|>{\raggedright\arraybackslash}p{8cm}|c|}
\hline
\textbf{№} & \textbf{Пошуковий рядок (TITLE-ABS-KEY)} & \textbf{Результатів} \\
\hline
1 & \texttt{"cooperative learning" AND "primary school" AND "cognitive activity"} & 28 \\
\hline
2 & \texttt{"collaborative learning" AND "elementary school" AND "inquiry"} & 41 \\
\hline
3 & \texttt{"group interaction" AND "primary education" AND "cognitive development"} & 17 \\
\hline
4 & \texttt{"inquiry-based learning" AND "primary school" AND "science"} & 54 \\
\hline
5 & \texttt{"project-based learning" AND "elementary school" AND "collaborative"} & 38 \\
\hline
6 & \texttt{"scientific inquiry" AND "young learners" AND "group work"} & 12 \\
\hline
7 & \texttt{"integrated curriculum" AND "primary school" AND "social studies" AND "science"} & 22 \\
\hline
8 & \texttt{"cross-curricular" AND "elementary education" AND "inquiry" AND "group"} & 9 \\
\hline
9 & \texttt{"STEAM" AND "primary school" AND "collaborative learning"} & 19 \\
\hline
10 & \texttt{"teacher facilitation" AND "cooperative learning" AND "primary"} & 14 \\
\hline
11 & \texttt{"scaffolding" AND "group work" AND "cognitive activity" AND "elementary"} & 11 \\
\hline
12 & \texttt{"formative assessment" AND "cooperative learning" AND "primary school"} & 8 \\
\hline
13 & \texttt{"peer assessment" AND "group interaction" AND "elementary education"} & 7 \\
\hline
\multicolumn{2}{|r|}{\textbf{Усього (з урахуванням дублікатів):}} & \textbf{280} \\
\hline
\multicolumn{2}{|r|}{\textbf{Після відбору за критеріями релевантності:}} & \textbf{63} \\
\hline
\multicolumn{2}{|r|}{\textbf{Безпосередньо процитовано:}} & \textbf{46} \\
\hline
\end{longtable}
\end{center}

\textbf{Критерії включення:} емпіричні або оглядові публікації; учні початкової школи (6--11 років); групова, кооперативна або колаборативна взаємодія у пізнавальній/дослідницькій діяльності; 2015--2026 рр.; мова~--- англійська або українська; тип~--- article, review.

\textbf{Критерії виключення:} виключно дошкільний або середній/старший шкільний рівень; вища освіта; роботи без зв'язку з класно-урочним контекстом; конференційні матеріали обсягом менше 6 сторінок без рецензування.

\newpage

% ============================================================
% ДОДАТОК Б — СПИСОК СКОРОЧЕНЬ І ТЕРМІНІВ
% ============================================================
\section*{Додаток Б}
\addcontentsline{toc}{section}{Додаток Б. Список скорочень та глосарій ключових понять}

\begin{center}
\textbf{Список скорочень}
\end{center}

\begin{description}[leftmargin=3cm, style=nextline]
  \item[ЯДС] «Я досліджую світ»~--- інтегрований курс для 1--4 класів початкової школи НУШ.
  \item[НУШ] Нова українська школа~--- концепція реформування загальної середньої освіти в Україні (2016~-- дотепер).
  \item[IBL] Inquiry-based learning~--- навчання на основі дослідження / дослідницьке навчання.
  \item[PjBL] Project-based learning~--- навчання на основі проєктів / проєктне навчання.
  \item[STEM] Science, Technology, Engineering, Mathematics~--- природничі науки, технології, інженерія, математика.
  \item[STEAM] STEM + Arts~--- STEM із включенням мистецтва та гуманітарних дисциплін.
  \item[ECEE] Explore–Create–Express–Evaluate~--- педагогічна рамка «Досліджуй–Створюй–Виражай–Оцінюй».
  \item[TIMSS] Trends in International Mathematics and Science Study~--- міжнародне дослідження досягнень учнів з математики та природничих наук.
  \item[ICT] Information and Communication Technologies~--- інформаційно-комунікаційні технології.
\end{description}

\vspace{1em}
\begin{center}
\textbf{Глосарій ключових понять}
\end{center}

\textbf{Кооперативне навчання} (cooperative learning)~--- педагогічна стратегія, в якій учні працюють у малих структурованих групах із позитивною взаємозалежністю, індивідуальною відповідальністю та чітким розподілом ролей; орієнтована на досягнення спільного навчального результату.

\textbf{Колаборативне навчання} (collaborative learning)~--- гнучка форма спільної навчальної діяльності, де учасники на рівних правах спільно конструюють нові знання та смисли без жорсткого розподілу функцій; акцент~--- на процесі взаємодії.

\textbf{Дослідницьке навчання} (inquiry-based learning, IBL)~--- підхід, заснований на залученні учнів до наукових практик: спостереження, формулювання запитань, висунення і перевірки гіпотез, аналізу даних та побудови пояснень. Виокремлюють структуроване, спрямоване і відкрите IBL залежно від ступеня самостійності учнів.

\textbf{Фасилітація} (facilitation)~--- роль педагога як організатора і підтримувача процесу спільного навчання, а не транслятора готових знань. Передбачає використання відкритих запитань, підтримку учнівських ідей і управління груповою динамікою.

\textbf{Фейсворк} (facework)~--- підтримання соціального іміджу та гідності учасників взаємодії (за Е.~Гоффманом); у контексті групового навчання~--- комунікативні стратегії, що запобігають публічному приниженню і зберігають психологічну безпеку в групі.

\textbf{Скафолдинг} (scaffolding)~--- педагогічна підтримка учнів на тому рівні, що відповідає зоні найближчого розвитку (за Л.\,С.~Виготським); поступово знімається в міру зростання самостійності учня.

\end{document}
